% ─────────────────────────────────────────────────────────────────────────────
\section{Introduction}
\label{sec:intro}
% ─────────────────────────────────────────────────────────────────────────────

This Low-Level Design (LLD) document describes the complete class-level architecture of
\textbf{MoonAI}: a predator-prey simulation for studying NeuroEvolution of Augmenting
Topologies (NEAT) and evolutionary algorithms.  It serves as the authoritative reference
linking the High-Level Design (HLD, December 2025) to the actual C++17 implementation.

The HLD defined six subsystems at a conceptual level.  This document closes the gap to the
source code: every class, struct, enum, algorithm, and design decision is documented with
enough precision that a developer could re-implement any component independently.

% ─────────────────────────────────────────────────────────────────────────────
\subsection{Object Design Trade-offs}
\label{sec:tradeoffs}
% ─────────────────────────────────────────────────────────────────────────────

Five key trade-offs shaped the implementation.

\subsubsection{Value Semantics vs.\ Pointer Semantics for \texttt{Genome}}

\texttt{EvolutionManager::population\_} stores \texttt{Genome} objects by value
(\texttt{std::vector<Genome>}).  \texttt{Species::members\_} holds raw, non-owning
\texttt{Genome*} pointers into that contiguous vector.

\textbf{Alternatives considered:}
\begin{itemize}[nosep]
  \item \texttt{std::shared\_ptr<Genome>} — would eliminate the raw pointer but incurs
    reference-counting overhead and heap fragmentation for every genome in the hot
    evaluate/reproduce loop.
  \item \texttt{std::unique\_ptr<Genome>} — moves ownership but breaks the contiguous
    layout needed for GPU packing.
\end{itemize}

\textbf{Decision:} The species list is rebuilt at the start of every generation
\emph{before} the population vector is reallocated, so the raw pointers are valid for the
entire lifetime of one generation.  The value-semantics vector gives cache-friendly
iteration and trivial move/copy semantics for checkpointing.

\subsubsection{GPU CSR Layout vs.\ Padded Layout}

The GPU inference path packs all agent networks into flat device arrays using Compressed
Sparse Row (CSR) encoding.  Each agent is described by a \texttt{GpuNetDesc} containing
eight integer offsets into shared flat arrays.

\textbf{Alternatives considered:}
\begin{itemize}[nosep]
  \item \textbf{Padded layout} — each network occupies a fixed-size slot
    (\texttt{MAX\_NODES} × \texttt{MAX\_NODES}).  Simple indexing, but wastes up to
    10× VRAM for the small early-generation topologies that NEAT produces.
  \item \textbf{Per-agent device allocation} — flexible but requires thousands of
    \texttt{cudaMalloc} calls and non-coalesced memory access.
\end{itemize}

\textbf{Decision:} CSR eliminates VRAM waste and enables coalesced reads along the
\texttt{eval\_order} and \texttt{conn\_w} arrays.  The host-side packing step
(\texttt{upload\_network\_data}) runs once per generation, so its O(N) complexity is
acceptable.

\subsubsection{Precomputed Incoming Adjacency in \texttt{NeuralNetwork}}

At construction, \texttt{NeuralNetwork} builds \texttt{incoming\_[node\_idx]}, a list of
\texttt{(source\_idx, weight)} pairs for each node.  The \texttt{activate()} method then
iterates this precomputed structure.

\textbf{Alternative:} Build the adjacency on each forward pass from the raw connection
list.  Profiling showed this as the dominant cost in the CPU inference path (map
allocation in the hot loop).

\textbf{Decision:} One-time O(E) build cost at network construction; zero allocation per
\texttt{activate()} call.  Activation state is stored in a pre-allocated \texttt{float}
vector \texttt{values\_} of fixed size.

\subsubsection{Static Utility Classes (\texttt{Mutation}, \texttt{Crossover},
\texttt{Physics})}

These classes expose only \texttt{static} methods and carry no instance state.

\textbf{Alternative:} Instantiable objects or singletons.  Singletons would require
locking for potential future multi-threading; instantiable objects would add unnecessary
storage.

\textbf{Decision:} Pure functions are simpler to reason about, easier to test (no
construction overhead), and completely thread-safe given the same \texttt{Random}
instance.  All operators are deterministic given the same \texttt{Random} state, so unit
tests verify correctness via outcome assertions rather than mock objects.

\subsubsection{\texttt{TickCallback} Decoupling}

\texttt{EvolutionManager} exposes a
\texttt{std::function<void(int, const SimulationManager\&)>} type alias called
\texttt{TickCallback} rather than holding a raw \texttt{Logger*}.

\textbf{Alternative:} Pass a \texttt{Logger*} directly — simpler but couples the
evolution engine to the data layer.

\textbf{Decision:} The callback pattern lets \texttt{main.cpp} wire logging, live
visualization updates, and benchmarking probes independently.  The slight
\texttt{std::function} overhead (one indirect call per tick) is negligible compared with
the simulation physics cost.

% ─────────────────────────────────────────────────────────────────────────────
\subsection{Interface Documentation Guidelines}
\label{sec:guidelines}
% ─────────────────────────────────────────────────────────────────────────────

\begin{itemize}
  \item All public interfaces are defined exclusively in \texttt{.hpp} files under
    \texttt{src/}.
  \item Include paths are written relative to the \texttt{src/} root, e.g.\
    \texttt{\#include "evolution/genome.hpp"}.
  \item \textbf{\texttt{const} correctness:} all read-only accessors are
    \texttt{const}; mutation methods are clearly non-const.
  \item \textbf{RAII:} all resources are managed by constructors/destructors —
    GPU memory (\texttt{GpuBatch}), file handles (\texttt{Logger}), and the SFML
    window (\texttt{VisualizationManager}).
  \item All symbols reside in namespace \texttt{moonai}; CUDA internals reside in
    \texttt{moonai::gpu}.
  \item This document uses exact C++17 signatures as the canonical interface specification.
\end{itemize}

% ─────────────────────────────────────────────────────────────────────────────
\subsection{Engineering Standards}
\label{sec:standards}
% ─────────────────────────────────────────────────────────────────────────────

\begin{itemize}
  \item \textbf{UML 2.5} — Class diagrams use standard notation: \texttt{+} public,
    \texttt{-} private, \texttt{\#} protected; \texttt{<<abstract>>},
    \texttt{<<static>>} stereotypes; standard relationship arrows (inheritance,
    composition, dependency).
  \item \textbf{IEEE 1016-2009} — \emph{Software Design Descriptions}: the structure of
    this document follows this standard.
  \item \textbf{C++17 ISO/IEC 14882:2017} — implementation language.
  \item \textbf{NEAT specification} — Stanley \& Miikkulainen (2002)~\cite{stanley2002}
    defines correctness requirements for the compatibility-distance formula, speciation,
    and reproduction.
  \item \textbf{ACM Code of Ethics} — reproducibility and transparency are maintained
    via seeded RNG, checkpointing, and open source tooling.
\end{itemize}

% ─────────────────────────────────────────────────────────────────────────────
\subsection{Definitions, Acronyms, and Abbreviations}
\label{sec:defs}
% ─────────────────────────────────────────────────────────────────────────────

See Section~\ref{sec:glossary} for the full glossary.  Key terms used throughout this
document:

\begin{description}[style=nextline,leftmargin=3.5cm,font=\ttfamily]
  \item[NEAT]  NeuroEvolution of Augmenting Topologies — the evolutionary algorithm
    implemented in \texttt{moonai::evolution}.
  \item[CSR]   Compressed Sparse Row — the GPU memory layout for variable-topology
    networks.
  \item[RAII]  Resource Acquisition Is Initialization — C++ idiom for deterministic
    resource lifetime.
  \item[RNG]   Random Number Generator; the project uses MT19937-64 via
    \texttt{std::mt19937\_64}.
  \item[NN]    Neural Network — a \texttt{NeuralNetwork} object built from a
    \texttt{Genome}.
  \item[HLD]   High-Level Design document (December 2025).
  \item[LLD]   This document — Low-Level Design.
  \item[SensorInput]  The 15-element observation vector fed to each agent's NN per tick.
  \item[TickCallback] \texttt{std::function<void(int, const SimulationManager\&)>} —
    hook invoked after each simulation tick.
  \item[GpuBatch] RAII wrapper for all CUDA device memory needed for one generation.
  \item[InnovationTracker] Per-generation map from connection pairs to innovation
    numbers, ensuring history alignment for NEAT crossover.
\end{description}
